
\begin{center}\Large\bfseries Дополнительные материалы к статье о СМО с отрицательными заявками\end{center}
\section*{1. Назначение и исходные данные}
Приложение сопровождает версию 9 статьи «Расчёт многоканальных СМО с отрицательными заявками и повторным обслуживанием». Оно содержит сверку метода Такахаси--Таками, погрешность двухфазного закона во всех режимах, поправку квадратуры для сохранённой работы, прежние серии имитационного моделирования и расширенный эксперимент. Везде ДИ означает полуширину 95\%-го доверительного интервала по распределению Стьюдента, а \(\Delta=100(\text{расчёт}-\text{ИМ})/\text{ИМ}\).

\section*{2. Метод Такахаси--Таками}
Стационарные вероятности двухфазной модели рассчитывались итерациями Такахаси--Таками в комплексной арифметике (библиотека Most-Queue, допуск \(10^{-12}\)), моменты --- по уравнениям меченой заявки. Те же уравнения решены прямым разреженным методом (таблица~XREFX0X). При \(\mathrm{CV}=0{,}8\) параметры \(H_2\) вещественны, но один вес отрицателен; при \(\mathrm{CV}=0{,}6\) и \(0{,}5\) параметры комплексно-сопряжённые. Мнимые части итоговых характеристик не превышают \(10^{-7}\).

\noindent TABLECAPX Таблица SEQXtab_1X. Сверка метода Такахаси--Таками с прямым разреженным решением тех же уравнений: \(n=3\), \(\lambda=1\), \(\delta=0{,}3\), \(N=200\), \(H=60\). Разность --- максимум по \(w_1,w_2,v_1,v_2\).\par

\setlength{\tabcolsep}{2pt}
\begin{tabular}{p{1.2cm}p{4.2cm}p{2.6cm}p{3.6cm}p{3.6cm}}
\toprule
$\mathrm{CV}$ & Механизм & Параметры $H_2$ & Отн. разность моментов & Невязка баланса \\
\midrule
1,2 & RCS, удаление & положит. & \(1{,}3\cdot10^{-11}\) & \(9{,}3\cdot10^{-14}\) \\
1,2 & Очищение & положит. & \(5{,}5\cdot10^{-12}\) & \(3{,}8\cdot10^{-13}\) \\
0,8 & RCS, удаление & отриц. вес & \(3{,}4\cdot10^{-10}\) & \(1{,}2\cdot10^{-12}\) \\
0,8 & Очищение & отриц. вес & \(5{,}6\cdot10^{-14}\) & \(1{,}1\cdot10^{-13}\) \\
0,6 & RCS, удаление & компл. & \(1{,}5\cdot10^{-7}\) & \(1{,}6\cdot10^{-9}\) \\
0,6 & Очищение & компл. & \(2{,}9\cdot10^{-13}\) & \(5{,}9\cdot10^{-13}\) \\
0,5 & RCS, удаление & компл. & \(1{,}0\cdot10^{-9}\) & \(2{,}9\cdot10^{-12}\) \\
0,5 & Очищение & компл. & \(6{,}8\cdot10^{-13}\) & \(1{,}2\cdot10^{-13}\) \\
\bottomrule
\end{tabular}




Библиотечный расчёт моментов численным дифференцированием ПЛС даёт то же среднее ожидание. Для пребывания при RCS он использует постоянную интенсивность \(\delta/m\), где \(m\) --- число занятых каналов в момент начала обслуживания; при \(\mathrm{CV}=1{,}2\) это завышает \(v_1\) с \(1{,}92203\) до \(1{,}92229\), при \(\mathrm{CV}=0{,}5\) --- с \(2{,}15436\) до \(2{,}15508\).

\noindent TABLECAPX Таблица SEQXtab_2X. Погрешность двухфазного закона во всех рассмотренных режимах относительно положительного фазового представления гамма-распределения с 16 узлами.\par

\setlength{\tabcolsep}{2pt}
\begin{tabular}{p{1.0cm}p{0.8cm}p{1.2cm}p{1.2cm}p{1.5cm}p{2.0cm}p{2.0cm}p{2.0cm}p{2.0cm}}
\toprule
$\mathrm{CV}$ & $n$ & $\rho$ & $\delta$ & Мех. & $w_1$ ($H_2$) & $\Delta_W$, \% & $v_1$ ($H_2$) & $\Delta_V$, \% \\
\midrule
0,8 & 3 & 0,70 & 0,30 & RCS & 0,4104 & 0,11 & 2,0850 & 0,016 \\
0,8 & 3 & 0,70 & 0,30 & Очищ. & 0,1925 & 0,22 & 1,4830 & −0,003 \\
0,8 & 3 & 0,84 & 0,30 & RCS & 0,8740 & 0,07 & 2,5878 & 0,021 \\
0,8 & 3 & 0,84 & 0,30 & Очищ. & 0,2962 & 0,20 & 1,5442 & 0,003 \\
0,8 & 2 & 0,63 & 0,15 & RCS & 0,6858 & 0,07 & 2,4831 & 0,016 \\
0,8 & 2 & 0,63 & 0,15 & Очищ. & 0,4433 & 0,11 & 1,9957 & 0,015 \\
0,6 & 3 & 0,70 & 0,30 & RCS & 0,4014 & −0,22 & 2,1362 & −0,032 \\
0,6 & 3 & 0,70 & 0,30 & Очищ. & 0,2011 & −0,48 & 1,5586 & 0,000 \\
0,6 & 3 & 0,84 & 0,30 & RCS & 0,8661 & −0,15 & 2,6365 & −0,042 \\
0,6 & 3 & 0,84 & 0,30 & Очищ. & 0,3112 & −0,42 & 1,6210 & −0,013 \\
0,6 & 2 & 0,63 & 0,15 & RCS & 0,6331 & −0,14 & 2,4744 & −0,030 \\
0,6 & 2 & 0,63 & 0,15 & Очищ. & 0,4283 & −0,23 & 2,0378 & −0,030 \\
0,5 & 3 & 0,70 & 0,30 & RCS & 0,3947 & −0,49 & 2,1544 & −0,070 \\
0,5 & 3 & 0,70 & 0,30 & Очищ. & 0,2039 & −1,10 & 1,5912 & −0,005 \\
0,5 & 3 & 0,84 & 0,30 & RCS & 0,8567 & −0,32 & 2,6503 & −0,090 \\
0,5 & 3 & 0,84 & 0,30 & Очищ. & 0,3165 & −0,97 & 1,6539 & −0,035 \\
0,5 & 2 & 0,63 & 0,15 & RCS & 0,6079 & −0,30 & 2,4672 & −0,065 \\
0,5 & 2 & 0,63 & 0,15 & Очищ. & 0,4199 & −0,52 & 2,0540 & −0,067 \\
\bottomrule
\end{tabular}




\section*{3. Сохранённая работа}
В основных примерах \(J=4\), \(N=100\); для RCS \(H=24\). Экстраполяция по \(1/K\) при \(K=8\), 12 и 16 имеет веса 2, \(-9\) и 8 и применяется к начальным моментам; в расширенном эксперименте \(J=3\), \(K=4\), 8 и 12 с весами \(1/2\), \(-4\) и \(9/2\). Таблица~XREFX1X содержит конечные сетки до поправки.

\noindent TABLECAPX Таблица SEQXtab_3X. Конечные сетки: вариант и \(\mathrm{CV}\), число фаз на класс, начальные моменты ожидания, время расчёта в секундах.\par

\setlength{\tabcolsep}{2pt}
\begin{tabular}{p{4.44cm}p{1.13cm}p{3.31cm}p{3.31cm}p{2.46cm}}
\toprule
Вариант, $\mathrm{CV}$ & $K$ & $w_1$ & $w_2$ & Время, с \\
\midrule
RCS, 0,8 & 4 & 1,239397 & 6,322977 & 7,6 \\
RCS, 0,8 & 8 & 1,139711 & 5,206458 & 120,1 \\
RCS, 0,8 & 12 & 1,105919 & 4,853345 & 83,0 \\
RCS, 0,8 & 16 & 1,088892 & 4,680347 & 337,2 \\
RCS, 1,2 & 4 & 1,864266 & 15,055481 & 8,1 \\
RCS, 1,2 & 8 & 1,726325 & 12,552254 & 124,3 \\
RCS, 1,2 & 12 & 1,679069 & 11,750098 & 126,2 \\
RCS, 1,2 & 16 & 1,655167 & 11,355380 & 520,7 \\
Все, 0,8 & 4 & 1,421020 & 8,207528 & 1,4 \\
Все, 0,8 & 8 & 1,347865 & 7,216742 & 40,5 \\
Все, 0,8 & 12 & 1,322678 & 6,894780 & 178,7 \\
Все, 0,8 & 16 & 1,309931 & 6,735760 & 1146,4 \\
Все, 1,2 & 4 & 2,324342 & 23,043622 & 10,3 \\
Все, 1,2 & 8 & 2,267271 & 21,460974 & 173,0 \\
Все, 1,2 & 12 & 2,247481 & 20,944201 & 115,2 \\
Все, 1,2 & 16 & 2,237492 & 20,690638 & 1357,7 \\
\bottomrule
\end{tabular}




Поправка квадратуры вычисляется по системе без воздействий на тех же сетках \(J,K\) и с тем же усечением \(N\) (таблица~XREFX2X). Эталон при \(\delta=0\) --- положительное фазовое представление гамма-распределения с 16 узлами или само \(H_2\). Поправка всегда положительна: квадратура с конечным \(J\) занижает время ожидания. Обоснованность поправки проверена по последовательности \(J=3,\ldots,6\) (таблица~XREFX3X): значения без поправки сходятся к пределу снизу, значения с поправкой практически от \(J\) не зависят.

\noindent TABLECAPX Таблица SEQXtab_4X. Поправка квадратуры для сохранённой работы. \(w_1(0)\) --- среднее ожидание при \(\delta=0\) на той же сетке (модель) и для исходного закона (эталон). Условия 9--12 --- из расширенного эксперимента.\par

\setlength{\tabcolsep}{2pt}
\begin{tabular}{p{4.4cm}p{2.2cm}p{2.2cm}p{2.2cm}p{1.9cm}p{2.2cm}}
\toprule
Вариант & $w_1$ без попр. & $w_1(0)$, модель & $w_1(0)$, эталон & Сдвиг, \% & $w_1$ с попр. \\
\midrule
RCS, 0,8 & 1,0373 & 0,9523 & 0,9537 & 0,143 & 1,0388 \\
Все каналы, 0,8 & 1,2711 & 0,9523 & 0,9537 & 0,116 & 1,2726 \\
RCS, 1,2 & 1,5824 & 1,3786 & 1,3851 & 0,410 & 1,5888 \\
Все каналы, 1,2 & 2,2071 & 1,3786 & 1,3851 & 0,294 & 2,2136 \\
Условие 9 & 0,7065 & 0,5805 & 0,5825 & 0,280 & 0,7085 \\
Условие 10 & 6,2171 & 4,1205 & 4,1368 & 0,262 & 6,2334 \\
Условие 11 & 0,8813 & 0,5805 & 0,5825 & 0,224 & 0,8833 \\
Условие 12 & 9,9564 & 2,7287 & 2,7437 & 0,150 & 9,9713 \\
\bottomrule
\end{tabular}




\noindent TABLECAPX Таблица SEQXtab_5X. Проверка поправки при \(n=2\), \(\rho=0{,}5\), \(\mathrm{CV}=0{,}8\), RCS, \(\delta=0{,}04\), \(K=8\), 12, 16. Предел последовательности без поправки по геометрической экстраполяции равен 0,70846.\par

\setlength{\tabcolsep}{2pt}
\begin{tabular}{p{1.5cm}p{4.0cm}p{4.0cm}}
\toprule
$J$ & $w_1$ без поправки & $w_1$ с поправкой \\
\midrule
3 & 0,70636 & 0,70849 \\
4 & 0,70757 & 0,70849 \\
5 & 0,70808 & 0,70848 \\
6 & 0,70830 & 0,70846 \\
\bottomrule
\end{tabular}




Основная таблица статьи использует серию из 32 прогонов длиной \(2\cdot10^6\) уходов. Прежняя серия приведена в таблице~XREFX4X; её интервалы примерно в четыре раза шире.

\noindent TABLECAPX Таблица SEQXtab_6X. Расчёт с поправкой и прежняя серия ИМ для сохранённой работы (8 прогонов, исключение \(2\cdot10^4\) и наблюдение \(2\cdot10^5\) уходов).\par

\setlength{\tabcolsep}{2pt}
\begin{tabular}{p{4.16cm}p{1.61cm}p{2.83cm}p{1.28cm}p{1.61cm}p{2.83cm}p{1.28cm}}
\toprule
 & \multicolumn{3}{c}{$w_1$} & \multicolumn{3}{c}{$v_1$} \\
Вариант & Расчёт & ИМ $\pm$ ДИ & $\Delta$, \% & Расчёт & ИМ $\pm$ ДИ & $\Delta$, \% \\
\midrule
RCS, 0,8 & 1,0388 & 1,0323 ± 0,0239 & 0,63 & 3,1695 & 3,1615 ± 0,0265 & 0,25 \\
Все каналы, 0,8 & 1,2726 & 1,2506 ± 0,0242 & 1,76 & 3,4473 & 3,4217 ± 0,0259 & 0,75 \\
RCS, 1,2 & 1,5888 & 1,6074 ± 0,0365 & −1,15 & 3,7344 & 3,7549 ± 0,0382 & −0,54 \\
Все каналы, 1,2 & 2,2136 & 2,2168 ± 0,0786 & −0,14 & 4,4274 & 4,4312 ± 0,0847 & −0,09 \\
\bottomrule
\end{tabular}




Сопоставление более грубой тройки фаз, числа классов и другого аргумента экстраполяции относится к результату без поправки:

\noindent TABLECAPX Таблица SEQXtab_7X. Относительные изменения, \%, относительно основной экстраполяции. Они не являются строгими границами ошибки.\par

\setlength{\tabcolsep}{2pt}
\begin{tabular}{p{4.44cm}p{2.79cm}p{2.79cm}p{2.79cm}p{2.79cm}}
\toprule
Вариант, $\mathrm{CV}$ & $w_1$: грубее & $w_1$: $J=3$ & $w_1$: $1/\sqrt K$ & $\sigma_W$: $1/\sqrt K$ \\
\midrule
RCS, 0,8 & 0,0191 & −0,1228 & −0,2201 & 1,0739 \\
RCS, 1,2 & 0,0177 & −0,4003 & −0,3029 & 0,8178 \\
Все, 0,8 & 0,0019 & −0,0854 & −0,2109 & 0,4408 \\
Все, 1,2 & −0,0176 & −0,2377 & −0,0752 & 0,3036 \\
\bottomrule
\end{tabular}




\section*{4. Расширенный эксперимент}
В каждом условии \(b_1=2{,}1\), \(\lambda=\rho n/b_1\). Исходный закон при \(\mathrm{CV}=0{,}8\) --- гамма-распределение, при \(\mathrm{CV}=1{,}2\) --- положительный \(H_2\). Удаление и очищение рассчитаны методом Такахаси--Таками, повторы --- совместной фазовой моделью (9 положительных фаз при \(\mathrm{CV}=0{,}8\)), сохранённая работа --- с поправкой квадратуры. В условиях 1--7 выполнено по 8 прогонов с исключением первых 50\,000 и наблюдением 500\,000 уходов. Для условий 8--12 добавлены серии из 32 прогонов длиной \(2\cdot10^6\) уходов: в прежней серии отдельные оценки этих условий выходили за интервалы. Прежние результаты условий 8--12 сохранены в исходных данных. Все средние и стандартные отклонения находятся внутри соответствующих интервалов; максимальные абсолютные расхождения средних ожидания и пребывания равны 1,36\% и 0,68\%, стандартных отклонений --- 1,11\% и 0,92\%.

\noindent TABLECAPX Таблица SEQXtab_8X. Дополнительные условия, метод расчёта (ТТ --- Такахаси--Таками, СФМ --- совместная фазовая модель) и серия ИМ. Номер сохраняется во всех следующих таблицах.\par

\setlength{\tabcolsep}{2pt}
\begin{tabular}{p{0.9cm}p{5.2cm}p{0.7cm}p{1.2cm}p{1.0cm}p{1.2cm}p{2.5cm}p{2.2cm}}
\toprule
№ & Механизм & $n$ & $\rho$ & $\mathrm{CV}$ & $\delta$ & Расчёт & Серия ИМ \\
\midrule
1 & RCS, удаление & 2 & 0,45 & 0,8 & 0,30 & ТТ & 8×0,5 млн \\
2 & RCS, удаление & 3 & 0,85 & 1,2 & 0,60 & ТТ & 8×0,5 млн \\
3 & Очищение & 2 & 0,90 & 0,8 & 0,15 & ТТ & 8×0,5 млн \\
4 & Очищение & 4 & 1,20 & 1,2 & 0,30 & ТТ & 8×0,5 млн \\
5 & RCS, новая длительность & 2 & 0,50 & 0,8 & 0,60 & СФМ & 8×0,5 млн \\
6 & RCS, новая длительность & 4 & 0,85 & 1,2 & 0,60 & СФМ & 8×0,5 млн \\
7 & Все каналы, новая длительность & 2 & 0,40 & 0,8 & 0,60 & СФМ & 8×0,5 млн \\
8 & Все каналы, новая длительность & 3 & 0,85 & 1,2 & 0,60 & СФМ & 32×2 млн \\
9 & RCS, сохранённая длительность & 2 & 0,50 & 0,8 & 0,04 & СФМ+попр. & 32×2 млн \\
10 & RCS, сохранённая длительность & 3 & 0,85 & 1,2 & 0,04 & СФМ+попр. & 32×2 млн \\
11 & Все каналы, сохранённая длительность & 2 & 0,50 & 0,8 & 0,04 & СФМ+попр. & 32×2 млн \\
12 & Все каналы, сохранённая длительность & 3 & 0,80 & 1,2 & 0,04 & СФМ+попр. & 32×2 млн \\
\bottomrule
\end{tabular}




\noindent TABLECAPX Таблица SEQXtab_9X. Средние в дополнительных условиях.\par

\setlength{\tabcolsep}{2pt}
\begin{tabular}{p{1.5cm}p{1.9cm}p{3.4cm}p{1.4cm}p{1.9cm}p{3.4cm}p{1.4cm}}
\toprule
 & \multicolumn{3}{c}{$w_1$} & \multicolumn{3}{c}{$v_1$} \\
№ & Расчёт & ИМ $\pm$ ДИ & $\Delta$, \% & Расчёт & ИМ $\pm$ ДИ & $\Delta$, \% \\
\midrule
1 & 0,1912 & 0,1904 ± 0,0011 & 0,45 & 1,7005 & 1,6999 ± 0,0027 & 0,04 \\
2 & 0,3195 & 0,3191 ± 0,0048 & 0,13 & 1,5240 & 1,5243 ± 0,0071 & −0,02 \\
3 & 0,9141 & 0,9139 ± 0,0071 & 0,02 & 2,3491 & 2,3493 ± 0,0072 & −0,01 \\
4 & 0,4585 & 0,4607 ± 0,0074 & −0,46 & 1,4992 & 1,5025 ± 0,0086 & −0,22 \\
5 & 0,9254 & 0,9214 ± 0,0061 & 0,44 & 3,3072 & 3,3018 ± 0,0078 & 0,16 \\
6 & 1,6984 & 1,7218 ± 0,0254 & −1,36 & 3,6576 & 3,6827 ± 0,0262 & −0,68 \\
7 & 0,7021 & 0,7015 ± 0,0058 & 0,08 & 3,2342 & 3,2340 ± 0,0081 & 0,01 \\
8 & 1,2073 & 1,2065 ± 0,0024 & 0,07 & 2,9820 & 2,9812 ± 0,0026 & 0,03 \\
9 & 0,7085 & 0,7098 ± 0,0016 & −0,18 & 2,9057 & 2,9073 ± 0,0020 & −0,06 \\
10 & 6,2334 & 6,2426 ± 0,0412 & −0,15 & 8,4140 & 8,4236 ± 0,0418 & −0,11 \\
11 & 0,8833 & 0,8828 ± 0,0019 & 0,05 & 3,1377 & 3,1373 ± 0,0024 & 0,01 \\
12 & 9,9713 & 10,0354 ± 0,0732 & −0,64 & 12,3128 & 12,3774 ± 0,0736 & −0,52 \\
\bottomrule
\end{tabular}




\noindent TABLECAPX Таблица SEQXtab_10X. Стандартные отклонения в дополнительных условиях.\par

\setlength{\tabcolsep}{2pt}
\begin{tabular}{p{1.5cm}p{1.9cm}p{3.4cm}p{1.4cm}p{1.9cm}p{3.4cm}p{1.4cm}}
\toprule
 & \multicolumn{3}{c}{$\sigma_W$} & \multicolumn{3}{c}{$\sigma_V$} \\
№ & Расчёт & ИМ $\pm$ ДИ & $\Delta$, \% & Расчёт & ИМ $\pm$ ДИ & $\Delta$, \% \\
\midrule
1 & 0,6101 & 0,6083 ± 0,0034 & 0,30 & 1,4398 & 1,4379 ± 0,0025 & 0,13 \\
2 & 0,8413 & 0,8399 ± 0,0150 & 0,17 & 1,6618 & 1,6621 ± 0,0123 & −0,02 \\
3 & 1,5759 & 1,5787 ± 0,0158 & −0,18 & 2,0143 & 2,0151 ± 0,0129 & −0,04 \\
4 & 1,0022 & 1,0060 ± 0,0157 & −0,38 & 1,5797 & 1,5834 ± 0,0115 & −0,23 \\
5 & 1,8110 & 1,8031 ± 0,0118 & 0,44 & 2,7302 & 2,7232 ± 0,0100 & 0,26 \\
6 & 2,6597 & 2,6896 ± 0,0602 & −1,11 & 3,5004 & 3,5260 ± 0,0485 & −0,72 \\
7 & 1,6145 & 1,6189 ± 0,0159 & −0,27 & 2,7742 & 2,7770 ± 0,0111 & −0,10 \\
8 & 2,0527 & 2,0526 ± 0,0056 & 0,00 & 2,8229 & 2,8232 ± 0,0045 & −0,01 \\
9 & 1,5126 & 1,5161 ± 0,0035 & −0,23 & 2,4453 & 2,4477 ± 0,0025 & −0,10 \\
10 & 7,9278 & 7,9450 ± 0,0660 & −0,22 & 8,4051 & 8,4221 ± 0,0626 & −0,20 \\
11 & 1,9296 & 1,9310 ± 0,0061 & −0,07 & 2,8428 & 2,8442 ± 0,0049 & −0,05 \\
12 & 13,1225 & 13,2506 ± 0,1305 & −0,97 & 13,5666 & 13,6922 ± 0,1271 & −0,92 \\
\bottomrule
\end{tabular}




\noindent TABLECAPX Таблица SEQXtab_11X. Стоимость исходных сеток расширенного эксперимента и максимальное абсолютное изменение начальных моментов при проверке отсечений. Для гамма-распределения время положительного PH-уточнения сохранено отдельно в исходных данных.\par

\setlength{\tabcolsep}{2pt}
\begin{tabular}{p{0.94cm}p{4.25cm}p{4.72cm}p{5.67cm}}
\toprule
№ & Время сеток, с & Пиковая память, МиБ & Изменение при уточнении \\
\midrule
1 & 1,8 & 95,1 & 0,000000 \\
2 & 1,0 & 107,8 & 0,000000 \\
3 & 1,7 & 95,9 & 0,000000 \\
4 & 1,2 & 115,4 & 0,000000 \\
5 & 1,3 & 101,0 & 0,000000 \\
6 & 1,5 & 116,2 & 0,000000 \\
7 & 1,3 & 101,8 & 0,000000 \\
8 & 1,0 & 107,7 & 0,000000 \\
9 & 8,6 & 331,5 & — \\
10 & 1579,9 & 367,9 & 0,001762 \\
11 & 6,4 & 99,5 & — \\
12 & 791,0 & 333,4 & 0,006041 \\
\bottomrule
\end{tabular}




Память указана по максимальному объёму оперативной памяти, занятой отдельным процессом (RSS); время --- по полному времени выполнения расчёта при одновременной работе нескольких процессов. Прочерк означает, что отдельное увеличение предела для данного условия не выполнялось.

\section*{5. Положительные фазовые распределения}
Для гамма-распределения при \(\mathrm{CV}=0{,}8\) использованы квадратуры с \(J=2\), 4 и 8, соответствующие 3, 5 и 9 фазам. На контрольной сетке аргумента максимальное отклонение функции распределения от исходного гамма-распределения составляет соответственно \(6{,}33\cdot10^{-3}\), \(8{,}08\cdot10^{-4}\) и \(9{,}54\cdot10^{-5}\). Это дискретная проверка, а не доказанная равномерная граница.

\noindent TABLECAPX Таблица SEQXtab_12X. Среднее ожидание при 3, 5 и 9 положительных фазах и изменение между 5 и 9 фазами, \%.\par

\setlength{\tabcolsep}{2pt}
\begin{tabular}{p{5.76cm}p{2.55cm}p{2.55cm}p{2.55cm}p{2.17cm}}
\toprule
Механизм & 3 фазы & 5 фаз & 9 фаз & $\Delta$, \% \\
\midrule
RCS, удаление & 0,409683 & 0,409965 & 0,409967 & 0,000550 \\
Очищение & 0,191780 & 0,192061 & 0,192062 & 0,000527 \\
RCS, новая длительность & 1,205480 & 1,205963 & 1,205966 & 0,000272 \\
Все каналы, новая длительность & 1,860551 & 1,859758 & 1,859769 & 0,000572 \\
\bottomrule
\end{tabular}




\noindent TABLECAPX Таблица SEQXtab_13X. Распределение Эрланга четвёртого порядка: \(\mathrm{CV}=0{,}5\), \(n=3\), \(\lambda=0{,}6\), \(\delta=0{,}3\), новая длительность. Отдельные 8 прогонов на вариант.\par

\setlength{\tabcolsep}{2pt}
\begin{tabular}{p{4.16cm}p{1.61cm}p{2.83cm}p{1.28cm}p{1.61cm}p{2.83cm}p{1.28cm}}
\toprule
 & \multicolumn{3}{c}{$w_1$} & \multicolumn{3}{c}{$v_1$} \\
Вариант & Расчёт & ИМ $\pm$ ДИ & $\Delta$, \% & Расчёт & ИМ $\pm$ ДИ & $\Delta$, \% \\
\midrule
RCS & 0,1992 & 0,1993 ± 0,0016 & −0,06 & 2,5824 & 2,5826 ± 0,0013 & −0,01 \\
Все каналы & 0,4335 & 0,4326 ± 0,0037 & 0,22 & 3,0838 & 3,0829 ± 0,0056 & 0,03 \\
\bottomrule
\end{tabular}




Двухфазный закон по трём моментам при \(\mathrm{CV}=0{,}8\) имеет вещественные параметры с отрицательным весом. Его ПЛС совпадает с ПЛС положительного распределения Кокса


\[\boldsymbol\alpha=(1,0),\]

\[S=\begin{pmatrix}-0{,}674119662&0{,}511518043\\0&-1{,}230642243\end{pmatrix},\]


с вероятностью перехода во вторую фазу 0{,}758794120. Поэтому при \(\mathrm{CV}=0{,}8\) расчёт методом Такахаси--Таками относится к вероятностному закону, а не только к алгебраическому приближению.

\section*{6. Воспроизводимость и границы проверки}
Исходные JSON-файлы содержат неокруглённые начальные моменты, параметры, отсечения, время и память, файлы прогонов --- начальные состояния генераторов случайных чисел. Длинные серии версии 9 используют базы 2026110810000--2026110816000 с шагом 1000 по вариантам и 8000, 9000 для условий 9 и 8 соответственно; номер прогона прибавляется к базе. Имитатор проверен при \(\delta=0\) на точном решении \(M/G/2\) с гамма-распределением.

Набор из 75 автоматических тестов охватывает известные частные случаи, сверку метода Такахаси--Таками с прямым решением и независимые реализации. Он не устанавливает равномерную точность для всех нагрузок и законов обслуживания. Разности сеток и поправка квадратуры являются численной диагностикой, а не доказанной оценкой погрешности.

